\scanpage{24}
\begin{center}-2.23-\end{center}
Therefore, in view of lemma 2.9, $\Delta(G)$ can be viewed as the result of pasting together spheres, one for each $T\subset G$, with sufficient identifications so that any two points are both on at least one sphere. The subsets $\Delta(T)$ will be called \uline{skeletons} of $\Delta(G)$. We can push this structure further in 2 ways: first the spheres $\Delta(T)$ can be broken up by considering the function $P(\delta)$:

\uline{Definition 2.7}: For all parabolic subgroups $P$, let $\Delta_P(G)$ be the set of $\delta\in\Delta(G)$ such that $P(\delta)\supset P$.

To describe $\Delta_P(G)$ explicitly, let $T$ be a maximal torus of $P$. Then $\Delta_P(G)\subset\Delta(T)$ by lemma 2.9. Suppose $P$ is spanned by $T$ and by additive subgroups $\mathbb G_a\subset P$ such that
\[
\alpha^{-1}\cdot\gamma\cdot\alpha=\chi(\alpha)\cdot\gamma
\]
for $\alpha\in T_k$, $\gamma\in(\mathbb G_a)_k$, and a suitable character $\chi$ of $T$. Then one checks that if $\delta\in\Delta(T)$ is induced by $\bar\lambda\in\unsure{\Gamma|\Pi}^{\mathbb R}(T)$, and if $\chi$ induces the linear functional $\bar\chi$ on $\unsure{\Gamma|\Pi}^{\mathbb R}(T)$, then
\[
P(\delta)\supset P\quad\Longleftrightarrow\quad\forall\text{ such }\chi,\quad\bar\chi(\bar\lambda)\leq0.
\]
Therefore, $\Delta_P(G)$ is simple the set of rational points on an inter-
section of hemi-spheres, (each determined by a rational linear functional, in fact). If $G$ is actually semi-simple, it can be shown that $\Delta_P(G)$ is even the set of rational points on a spherical simplex, and hence the collection of $\Delta_P(G)$ constitutes a ``triangulation'', and $\Delta(G)$ is a complex formed from the simplices $\Delta_P(G)$. This is the motivation
\mspage{23}
